
多核神经网络处理器（neural processing unit，NPU）的并行计算收益受到数据搬运与同步等待的制约。本文将计算图（computational graph）中的算子（operator）划分为子图（subgraph），联合确定子图的核心（core）归属与执行次序，在满足数据依赖和缓存容量约束的前提下，以总体完成时间（makespan）最小为主要目标，兼顾额外主存搬运量与求解耗时。针对三问逐步放宽的数据复用条件，分别设计结构化候选构造、计算与搬运协同划分、共享缓存下的计算次序优化方法。以下平均加速比（average speedup）以官方单核方案为基准，对 100 张图逐例计算后取算术平均。

\textbf{针对问题一，}子图独立封装为任务（task），跨任务数据经双倍数据率内存（double data rate memory，DDR）中转。本文依据连通分量（connected component）与拓扑前沿（topological frontier）构造候选，利用\textbf{桥树分解}（bridge-tree decomposition）识别可整体迁移的独立分支，在核心间重新分配计算；再按拓扑高度（topological level）安排执行次序，避免分核后形成循环等待。候选先按完成时间选优，持平时再比较搬运量。\textbf{2、3、4、5 核心的平均加速比分别为 1.9516、2.7514、3.4786、4.0553}；5 核心下 100 张图的额外 DDR 搬运总量为\textbf{1,053,684,536 字节}。

\textbf{针对问题二，}同核子图合并为一个任务，在私有缓存（private cache）中复用中间数据。本文利用流水线（pipeline）空闲区间安排计算，以\textbf{超图连接度}（hypergraph connectivity）刻画同一数据向不同核心复制的代价；通过带工作量约束的\textbf{最小割}（minimum cut）调整分核，防止为减少通信而将计算过度集中于单核。复制量用于构造候选，最终按完整模拟的完成时间选优。\textbf{2、3、4、5 核心的平均加速比分别为 2.3167、3.2356、3.9711、4.5498}；5 核心下 100 张图的额外 DDR 搬运总量为\textbf{866,992,066 字节}。

\textbf{针对问题三，}引入共享只读缓存（shared read-only cache）。对满足结构条件的归约树（reduction tree），将子树内部存储峰值记为 $h_i$、完成后保留的输出量记为 $r_i$，证明在子树互不交错、内部开销固定的模型下，按\textbf{$h_i-r_i$ 降序处理}可使模型内存储峰值最小；该结论不直接保证整图完成时间缩短。因此，结合完成时间下界（lower bound）排除不可能改善的候选，至多评测 3 次，仅采纳完成时间更短的方案。\textbf{2、3、4、5 核心的平均加速比分别为 2.3120、3.2587、4.0915、4.7576}，逐图字节命中率（byte hit rate）的算术平均分别为\textbf{16.25\%、25.96\%、24.53\%、31.36\%}。5 核心下 100 张图的额外 DDR 搬运总量为\textbf{1,325,640,272 字节}，该官方指标不扣除缓存命中的字节；固定方案与其余配置，开启缓存相较关闭缓存的平均成对加速比为\textbf{1.0829}。

三问覆盖 100 张图与 1 至 5 核心，\textbf{1500 种组合均得到合法方案}。与各自前一完整版本相比，三问分别在 500 种组合中的\textbf{48、268、24 种组合}上缩短完成时间，其余持平。共享主机上，三问的平均端到端求解耗时依次为\textbf{2.825 秒、4.029 秒和 4.401 秒}，表明结构化构造与有限评测能够在本组基准上兼顾方案质量和求解效率。完成时间由官方离散事件仿真（discrete-event simulation）给出，求解耗时按墙钟时间（wall-clock time）记录，均不代表真实芯片实测收益。
\keywords{计算图切分；多核调度；超图最小割；树状存储排序；共享只读缓存}
